\scanpage{168}
% L1–2
n'est autre que l'image inverse de cette
section.

% L3–12
\deleted{\unsure{L|On}}On voit que la bonne notion ici
est celle de fonctions \unsure{semi|sous} quadratiques
d'un torseur $A$ dans un autre
$Q$, sous $L^{\otimes 2}$ cette fois, ``compatible''
avec \deleted{\unsure{$\sigma$}} ($u\mathrel{\vdash\!\!-\!\!-}u^{\unsure{2}}$ de $L$ dans $L^{\otimes 2}$ et) $t$: ce
sont les fonctions \inserted{\unsure{para|sous}} quadratiques
\deleted{\unsure{$q$}} \deleted{telles}.
\[
q:A\longrightarrow Q
\]
telles que
\[
q(x+\underset{\substack{\rotatebox{-90}{$\in$}\\L}}{\alpha})=q(x)-
\underset{\substack{\rotatebox{-90}{$\in$}\\L}}{(t\wedge x)}\otimes\alpha+
\alpha\deleted{\unsure{\otimes}}\otimes\alpha
\]
% slanted
\margin{attention conv. de signe!}

% L13–17
Deux telles fonctions \inserted{\unsure{para|sous}} quadratiques diffèrent
par une constante. Si on en a une,
cela fixe $Q$ à iso. près comme
sol. du pb universel. Et on vient
d'en construire une tantôt!

% L18–20
\deleted{On peut}
Ainsi, à \unsure{tte|tt} ``droite affine'' $A$ à involution
on a associé $Q(A,\sigma)=Q(A)$
et $A\longrightarrow Q(A)$ rev. de degré $2$
% archive sheet number: 167
